\documentclass[11pt]{article}
\usepackage[margin=1in]{geometry}
\usepackage[T1]{fontenc}
\usepackage{lmodern}
\usepackage{microtype}
\usepackage{amsmath,amssymb,amsthm}
\usepackage{booktabs}
\usepackage{array}
\usepackage{natbib}
\usepackage{hyperref}
\usepackage{xr-hyper}
\externaldocument[AD-]{09_adaptability}
\externaldocument[GE-]{01_law_of_graphic_equalisation_theory_referenced}
\hypersetup{colorlinks=true,linkcolor=blue,citecolor=blue,urlcolor=blue}
% GENERATED by tools/build_paper_metadata.py from paper_metadata.json.
% DO NOT EDIT: edit paper_metadata.json, then regenerate this file.
\providecommand{\PaperDraft}{draft-2026-09-06.1}
\providecommand{\PaperDisplayDate}{6 September 2026}
\expandafter\def\csname PaperRevision@00_is\endcsname{129}
\expandafter\def\csname PaperRevision@01_law_of_graphic_equalisation_theory_referenced\endcsname{294}
\expandafter\def\csname PaperRevision@01a_graphic_equalisation_physics_grounding\endcsname{47}
\expandafter\def\csname PaperRevision@01b_graphic_equalisation_mind_and_self_grounding\endcsname{47}
\expandafter\def\csname PaperRevision@01c_graphic_equalisation_cosmology_grounding\endcsname{11}
\expandafter\def\csname PaperRevision@01d_graphic_equalisation_born_recovery\endcsname{6}
\expandafter\def\csname PaperRevision@01e_graphic_equalisation_communication_grounding\endcsname{6}
\expandafter\def\csname PaperRevision@01f_graphic_equalisation_qualia_and_consciousness_grounding\endcsname{3}
\expandafter\def\csname PaperRevision@01g_graphic_equalisation_total_perspective_vortex_grounding\endcsname{4}
\expandafter\def\csname PaperRevision@01h_graphic_equalisation_information_gravity_quantum_reconciliation\endcsname{5}
\expandafter\def\csname PaperRevision@01i_graphic_equalisation_passive_continuation_momentum_gravity\endcsname{12}
\expandafter\def\csname PaperRevision@01s_graphic_equalisation_statement\endcsname{294}
\expandafter\def\csname PaperRevision@02_seldon_manifolds_theory_referenced\endcsname{31}
\expandafter\def\csname PaperRevision@02a_seldon_manifolds_psychological_grounding\endcsname{15}
\expandafter\def\csname PaperRevision@02b_seldon_manifolds_social_and_psychological_application\endcsname{7}
\expandafter\def\csname PaperRevision@02c_seldon_manifolds_decision_making\endcsname{11}
\expandafter\def\csname PaperRevision@03_law_of_narrative_momentum_theory_referenced\endcsname{32}
\expandafter\def\csname PaperRevision@04_narrativium_theory_referenced\endcsname{23}
\expandafter\def\csname PaperRevision@05_frames_relation_and_the_empty_limit\endcsname{11}
\expandafter\def\csname PaperRevision@06_assessing_an_account_of_this_kind\endcsname{29}
\expandafter\def\csname PaperRevision@07_causal_attitude_networks_on_the_substrate\endcsname{4}
\expandafter\def\csname PaperRevision@08_certified_answers_as_an_audit\endcsname{4}
\expandafter\def\csname PaperRevision@09_adaptability\endcsname{94}
\expandafter\def\csname PaperRevision@09a_adaptability_implementation\endcsname{7}
\expandafter\def\csname PaperRevision@10_applied_graphic_equalisation\endcsname{7}
\expandafter\def\csname PaperRevision@10a_red_queen_race_civilisational_rate_of_change_llms_biological_timescales_anthropocene_transition\endcsname{13}
\expandafter\def\csname PaperRevision@11_relational_information_ethics\endcsname{16}
\expandafter\def\csname PaperRevision@A1_constructive_graphic_equalisation_implementation_referenced\endcsname{8}
\expandafter\def\csname PaperGenerated@01s_graphic_equalisation_statement\endcsname{1}
\edef\PaperMetadataStem{\jobname}
\expandafter\ifx\csname PaperRevision@\PaperMetadataStem\endcsname\relax
  \PackageError{paper_metadata}{No metadata for \PaperMetadataStem}{Add the paper stem to paper_metadata.json.}
  \def\PaperRevision{?}
\else
  \expandafter\let\expandafter\PaperRevision\csname PaperRevision@\PaperMetadataStem\endcsname
\fi
\newcommand{\PaperDateLine}{\textbf{Draft \texorpdfstring{\textemdash}{---}\ \PaperDraft}\\[2pt]\PaperDisplayDate\ \texorpdfstring{\textbullet}{\textperiodcentered}\ document revision \PaperRevision}
\newcommand{\PaperStatementDateLine}{\textbf{Draft \texorpdfstring{\textemdash}{---}\ \PaperDraft}\\[2pt]\PaperDisplayDate\ \texorpdfstring{\textbullet}{\textperiodcentered}\ statement}
\expandafter\ifx\csname PaperGenerated@\PaperMetadataStem\endcsname\relax
  \newcommand{\PaperFooterStamp}{\textsc{\PaperDraft\,r\PaperRevision}}
\else
  \newcommand{\PaperFooterStamp}{\textsc{\PaperDraft\ / statement / from long r\PaperRevision}}
\fi
\newcommand{\PaperNoticeStamp}{\textbf{DRAFT --- \PaperDraft\ r\PaperRevision} (\PaperDisplayDate).}


\newtheorem{definition}{Definition}
\newtheorem{proposition}[definition]{Proposition}
\newtheorem{remark}[definition]{Remark}

\title{Cost-Floor Compression, the AI Wedge, and Civilisational Transition\\\large A Conditional GE Assessment}
\author{Matthew Parslow\\Independent Researcher}
\date{\textbf{Draft \texorpdfstring{\textemdash}{---}\ \PaperDraft}\\[2pt]\PaperDisplayDate\ \texorpdfstring{\textbullet}{\textperiodcentered}\ document revision \PaperRevision}

\begin{document}
\maketitle
\begin{center}
\textit{Unreleased application draft. This paper supplies a measurement contract and conditional accounting results; it does not forecast civilisational collapse, establish a realised GE substrate, or recommend a political programme.}
\end{center}
\vspace{1em}

\begin{abstract}
AI is a visible leading edge of a broader question: when the audited cost of
performing specified transformations falls toward a measured feasible lower
envelope, can the institutional arrangements that allocate work, income,
access, and authority remain unchanged? This paper does not infer an answer
from a benchmark, an API tariff, agricultural productivity, or a logistics
platform. It defines an outcome-scaled cost ledger separating information,
coordination, embodied, verification, governance, and externality costs. A
nominally cheaper task is not thereby a cheaper successful social outcome.

The paper then states two conditional results. First, a fall in technical
capability cost alone does not identify delivered adaptation, equitable access,
or net benefit: it may also increase deployment, induced demand, verification
load, and residual harm. Second, given a declared GE transition premise that
the present civilisational regime cannot reproduce all of its defining
arrangements under measured cost-floor pressure, the relevant outcome is a
partition, not a prediction: either minimum viable reproduction fails
(terminal discontinuity), or viable activity continues through a successor
regime. The latter has an analytically open, not automatically benign, space of
possibilities. LLMs are treated as a cost and deployment variable, not as a
claim about agency or general intelligence. A hypothetical graph-based GE
substrate is an unverified additional scenario, not evidence. The paper states
data requirements, counterweights, and falsifiers for every empirical step.
\end{abstract}

\section{Scope: a conditional transition question}\label{sec:scope}

This paper supplies a conditional classification. If a named regime cannot
continue to reproduce its defining arrangements while the cost basis of its
principal transformations changes, then unchanged persistence is not the
relevant comparison class. The possible continuations are classified as either
failure to reproduce a declared viability vector or reproduction through a
materially different regime. Whether the antecedent holds, when it might hold,
and what a successor would be remain open empirical and political questions.

This application extends neither the physical claims of GE nor its proposed
graph-based substrate. Within the GE programme, narrative and narrativium name
a way to ask whether a collective arrangement can retain its identity while its
relations, costs, and means of reproduction are reorganised. They are not
substitutes for labour, productivity, price, energy, environmental, or
institutional data. A GE-derived implementation efficiency gain has not been
shown here. It is only a separately labelled scenario in
Section~\ref{sec:ge-scenario}.

\subsection{Analysis lens: current manifold and possibility set}

For this paper, the \emph{current manifold} $\mathcal{M}(R_0)$ is a declared
analytical state space of the present regime's reproduction arrangements,
ledger variables, and measured constraints. ``Manifold'' is a classificatory
term here, not a claim that the regime is a differentiable physical manifold.
Given $G(R_0)$, the analysis maps a possibility set containing terminal
discontinuity and viable non-isomorphic successor regimes. It assigns no value
ordering to those members. Its output is the declared state description,
conditions, and branch classification.

\subsection{Continuance interpretation}

A separate interpretive lens may then ask what the classification means for a
specified aim. This paper records the following continuance lens: civilisational
continuance---a trajectory that continues to satisfy a declared collective
viability condition---ranks above terminal discontinuity. Conditional on that
continuance, a trajectory ranks higher when it reduces measured human suffering
$S^{H}$ under a declared protocol and retains human participation $P^{H}$ in
the continuing system. Thus, under this lens, continuance is the best outcome;
within it, lower $S^{H}$ and continued $P^{H}$ define the ideal direction.

This ordering is an explicit interpretive input, not a consequence of GE, the
cost ledger, or the transition partition. It requires its own definitions of
civilisation, humanity, suffering, participation, population, horizon, and
measurement threshold. The analytical layer remains unchanged if a different
interpretive lens is supplied.

The empirical object is not ``civilisation'' at one unmeasurable speed. A first
test must name a unit $u$, outcome $o$, population, time interval, viability
threshold, and allocation boundary: for example, a public service's ability to
deliver a specified safe outcome, an agricultural supply system's ability to
meet a nutritional standard, or a jurisdiction's ability to maintain a
critical infrastructure service. A claim about civilisation requires an
explicit aggregation argument; this paper supplies none.

\section{Cost-floor pressure and the ledger}\label{sec:ledger}

\begin{definition}[Audited task cost]\label{def:audited-cost}
For named unit $u$, outcome $o$, task configuration $q$, and interval $i$, let
\[
 C^{\mathrm{aud}}_{i}(q)=C^{I}_{i}+C^{Q}_{i}+C^{B}_{i}+C^{V}_{i}+C^{G}_{i}+C^{E}_{i}
\]
be the cost required to produce a successful, accessible, and audited outcome.
The terms are respectively information/inference, coordination, embodied
material and energy, verification and repair, governance/authorisation, and
allocated externality costs. Each term requires a stated payer, location,
time horizon, allocation rule, and uncertainty range. A list API price is at
most one component of $C^{I}_{i}$.
\end{definition}

\begin{definition}[Feasible lower envelope and compression]\label{def:envelope}
Let $F_{u,o,i}$ be the lowest observed, boundary-harmonised audited cost
among configurations meeting fixed outcome, reliability, access and
externality conditions. Envelope compression means $F_{u,o,i+1}<F_{u,o,i}$
for declared intervals on the same outcome scale; report the rate
$(F_{u,o,i+1}-F_{u,o,i})/(t_{i+1}-t_i)$.
The separate gap $C_i^{\rm aud}(q)-F_{u,o,i}$ measures a configuration's
convergence to the frontier. It can rise when the frontier improves and
can be zero while all costs rise. Neither quantity is a thermodynamic floor.
\end{definition}
The empirical envelope differs from the logical-erasure floor discussed in
Graphic Equalisation's Remark~\ref{GE-rem:red-queen}. A declining observed
price does not compress Landauer's thermal bound, nor does it measure total
civilisational adaptation cost.

For a named partial ledger, take the AI Index's GPT-3.5-level benchmark
threshold, a million-token inference-price unit, its listed qualifying
configurations and the November 2022--October 2024 interval. The reported
price moves from $\$20$ to $\$0.07$ per million tokens \citep{stanford2025}.
This instantiates an inference-price frontier only: the operator's full
coordination, embodied, verification, governance and externality costs are
unmeasured. The price-compression claim fails if a re-audit of the same
threshold and token unit does not preserve the fall; the full audited-cost
compression claim is not established by this partial ledger.



\begin{remark}[Grammar-relative shapes, floors, and MDL]\label{rem:mdl-cost}
Here a declared outcome $o$ is a shape, or a shape trajectory, rather than an
invariant that a transformation must leave unchanged. A grammar $\gamma$ is a
declared repertoire of generative relations/operations; it makes a shape space
$\Sigma_{\gamma}$ available and induces a grammar-relative lower envelope
$F^{\gamma}_{u,o,i}$. A new grammar $\gamma'$ can change the available shape
space as well as the envelope. A grammar step is recorded through a declared
shape comparison map $\Pi_{\gamma\to\gamma'}$ when
\[
 F^{\gamma'}_{u,\Pi_{\gamma\to\gamma'}(o),i}
 < F^{\gamma}_{u,o,i}.
\]
It is thus a step change in what shape can be generated and at what audited
cost, not preservation of a fixed outcome invariant.

Minimum description length (MDL) selects a sufficient account by minimising
\[
 L(M)+L(D\mid M),
\]
the description of a model plus the residual description of its data
\citep{rissanen1978}. Under the GE reading used here, MDL pressure and
cost-floor pressure are the same compression pressure at different realisation
levels. MDL acts on an account and its generative grammar; the audited ledger
records the information, coordination, embodied, verification, authorisation,
and externality transformations through which shapes are realised.

This is not an equality $C^{\mathrm{aud}}=L(M)+L(D\mid M)$. The correspondence
requires a declared map from descriptive grammar to realised grammar, and every
cost component remains boundary-, substrate-, and frame-relative. In particular,
MDL alone does not establish an API price, energy use, total social cost, an
empirical lower envelope, a grammar step, or a civilisational transition.
Paper~01 treats the agreement between its constitutive MDL reading and the
account as a consistency condition rather than independent evidence; this
application retains that limit.
\end{remark}

The definition prevents a common equivocation. A system may reduce token price,
labour hours, or local cash expenditure while transferring work to users,
future maintenance, a less powerful population, an energy system, or an
environmental sink. Such a change has not established lower audited task cost.
Conversely, some costs can fall sharply without approaching a physical limit.
The relevant question is which declared component falls and what remains.

For a declared outcome, let $D_i$ be changed demand, $A_i$ verified delivered
adaptation, $H_i$ induced demand and residual harm, and $B_i$ recorded backlog:
\[
 B_{i+1}=\max\{0,B_i+\Delta D_i-A_i+H_i\}.\tag{1}
\]
Let $T_i$ denote task-relevant technical performance, $X_i$ deployment
exposure, $Q_i$ coordination quality, and $L_i$ response lag. Cost compression
is a possible input to delivered adaptation, but not its definition:
\[
 A_i=f(T_i,X_i,Q_i,L_i,C^{\mathrm{aud}}_i;\ \text{access, reliability, governance}).\tag{2}
\]
This notation allows a reduction in one cost term to help $A_i$, while the same
change increases $X_i$, $\Delta D_i$, or $H_i$.

\begin{proposition}[Cost is not delivered adaptation]\label{prop:cost-not-adaptation}
A measured decline in $C^{I}_i$, or in any other one component of
$C^{\mathrm{aud}}_i$, does not identify the sign of $\Delta D_i-A_i+H_i$.
\end{proposition}
\begin{proof}
Equation~(2) leaves access, reliability, governance, coordination, and the
other cost components unspecified. Hold the observed component fixed and
construct one admissible case in which lower cost enables verified service and
raises $A_i$, and another in which deployment grows but failures, fraud,
rework, energy use, or exclusion raise $H_i$ or $\Delta D_i$. Both cases share
the observed cost decline and give different signs in (1). Therefore the cost
decline alone identifies neither net adaptation nor backlog direction. \qed
\end{proof}

\begin{proposition}[Conditional backlog growth]\label{prop:backlog}
Fix a stable common outcome scale and intervals $i=1,\ldots,n$. If
\[
 \Delta D_i-A_i+H_i\geq\epsilon>0
\]
in every interval, with no unrecorded reserve, substitution, or threshold
redefinition, then $B_n\geq B_0+n\epsilon$.
\end{proposition}
\begin{proof}
The stated positive increment makes the lower clamp in (1) inactive as a lower
bound. Thus $B_{i+1}\geq B_i+\epsilon$; iteration proves the result. \qed
\end{proof}

A converse catch-up condition remains equally important. If, while $B_i>0$,
$A_i-\Delta D_i-H_i\geq\epsilon>0$, then
$B_{i+1}\leq\max\{0,B_i-\epsilon\}$; backlog falls by at least
$\min\{\epsilon,B_i\}$ in that interval. Cost compression cannot remove this
recovery possibility by rhetoric.

\section{What the evidence supports---and does not}\label{sec:evidence}

The available evidence supports changes in particular task costs and production
structures, not a universal cost-collapse law. The 2025 AI Index reports that
inference cost at a GPT-3.5-comparable performance level fell by more than a
factor of 280 from November 2022 to October 2024, alongside large advances on
named technical benchmarks \citep{stanford2025}. That is strong evidence for a
rapidly moving cost/performance frontier. It is not an all-in agent cost, a
reliability guarantee, a capability measure for collective judgement, or a
prediction of employment outcomes.

The agricultural record illustrates the broader, slower pattern without making
AI its cause. Agriculture is the clearest historical example of a primary
activity whose labour requirement and input composition have been transformed:
in historical series for France, the United Kingdom, Italy, and the Netherlands,
large agricultural employment shares are documented before industrialisation \citep{owidag2026}. This is country evidence, not a reconstructed
share for every human population or a single global historical estimate. USDA
accounting further reports that U.S. farm output nearly tripled from 1948--2023,
with aggregate inputs changing little; land and labour inputs fell while
equipment and intermediate inputs grew \citep{usda2026}. The appropriate
inference is labour, substitution, and productivity change under particular
historical conditions, not a reconstructed all-in economic-cost series,
costless food, or the disappearance of land, energy, water, biology, capital,
and distribution. For logistics, OECD analysis finds that
digital connectivity and trade-facilitation measures such as automated border
processes can be especially trade-enhancing for parcels, while regulatory and
physical frictions remain material \citep{oecd2021,oecd2025}. A platform's
convenience therefore cannot stand as evidence that commerce or transport is
near a universal floor.

The proposed evidence protocol is task-level. It requires repeated
performance-normalised cost data; the full audit boundary in Definition~\ref{def:audited-cost};
reliability, appeal, and repair rates; accessibility and distribution; energy,
water, material, and local infrastructure effects; and a comparison with an
unchanged or alternative configuration. A named LLM release is admissible only
with model/version, evaluation, prompt/tool setting, throughput, cache and
context assumptions, tariff time, and oversight. Announced or vendor-listed
pricing is a commercial input, not realised social cost.

\section{Economic reorganisation is a hypothesis to test}\label{sec:economy}

The civilisational concern is not that lower cost automatically makes people
poorer. It is that many current arrangements connect access, purchasing power,
and social status to paid work under conditions of costly production. If
cost-floor pressure affects a large enough set of tasks, the relationships
between labour income, ownership of productive assets, market demand, and
legitimacy become empirical variables rather than stable background conditions.

Regime existence and good governance are distinct. A regime can persist,
extract, coerce, exclude, or externalise harms without equitable access,
legitimacy, due process, or effective governance. None of those desirable
properties is an antecedent of the terminal/successor partition. They enter this
paper only when a named outcome claims them, or when their absence changes the
declared viability vector, reliability, cost boundary, or harm ledger. A
successor can therefore exist without being good; the partition does not
endorse it.

The ILO finds that roughly one quarter of global employment is potentially
exposed to generative AI, while judging task transformation more likely than
replacement in most cases \citep{ilo2023}. Its later evidence records a 1.6
percentage-point decline in the global labour-income share from 2004--2024 and
identifies several structural influences, including capital-augmenting
technology and worker bargaining power \citep{ilo2025}. Neither result proves
that AI causes a future income collapse. Together they justify measuring who
owns productive capacity, who receives gains, whether demand remains adequate
for the declared outcome, and whether institutions provide non-wage routes to
access and agency when those outcomes are at issue.

Competition and legitimacy are optional evaluative and pathway variables, not
requirements for mere system continuation. The UK CMA has identified risks to
fair, open, and effective competition in foundation-model markets \citep{cma2024}.
The IEA reports that social acceptability of energy-intensive data-centre
projects can depend on demonstrable local benefits and fair allocation of costs
\citep{iea2026}. These are evidence relevant to named equitable or accepted
deployment pathways, not conditions a system must meet in order to persist.

\begin{remark}[A testable economic transition claim]
For a named population, the economically relevant claim is not ``the economic
system is destroyed.'' It is: after a declared cost reduction, do verified
access, income security, agency, and outcome delivery keep pace with changes in
labour demand, asset ownership, prices, and environmental burden? A negative
result for one population or sector does not establish a civilisational result;
a positive result falsifies a claim of necessary local exclusion. A system can
continue despite either result; the result instead classifies the named quality
of continuation.
\end{remark}

\section{The GE-conditioned fork}\label{sec:fork}

\begin{definition}[Regime and transition premise]
Let $R_0$ be a declared present regime, represented by a set of defining
reproduction arrangements $\mathcal{I}_0$ and a minimum viability vector $v$.
The \emph{GE transition premise} $G(R_0)$ is the additional theoretical and
empirical premise that, under the stated cost-floor pressure, $R_0$ cannot
continue indefinitely while retaining all of $\mathcal{I}_0$. This paper does
not establish $G(R_0)$.

A \emph{terminal discontinuity} occurs when the declared viable reproduction
condition $v$ fails. A \emph{successor regime} occurs when viable reproduction
continues but the arrangement is no longer isomorphic to $\mathcal{I}_0$ under
a stated identity test. The successor may be more, less, or differently
inclusive, stable, sustainable, or free.
\end{definition}

\begin{proposition}[Conditional transition partition]\label{prop:fork}
Assume $G(R_0)$ and an exhaustive, mutually exclusive viability classification:
a post-transition trajectory either fails $v$ or continues to satisfy $v$.
Then every classified trajectory is either a terminal discontinuity or a
successor regime.
\end{proposition}
\begin{proof}
By the viability classification, a trajectory fails $v$ or satisfies $v$. In
the first case it meets the definition of terminal discontinuity. In the second,
$G(R_0)$ excludes indefinite retention of every defining arrangement in
$\mathcal{I}_0$, so it meets the definition of successor regime. The cases are
mutually exclusive by the classification. \qed
\end{proof}

This is a partition conditional on $G(R_0)$, not a date, probability, or
prediction. It gives no cardinal measure of the successor's possibility space.
It means only that a viable successor is not constrained to preserve the present
regime's full identity. A broad description such as ``the end of the
Anthropocene,'' ``phase change,'' ``singularity,'' or ``Lyapunov-horizon
collapse'' is admissible only as a declared analytical umbrella for this
question. It is not a formal geological epoch claim, demonstrated physical
singularity, known mechanism, or end of human influence or geology.

\section{Human timescales, physical constraints, and the GE scenario}\label{sec:ge-scenario}

Human biology is one constrained and socially mediated timescale among learning,
care, cohort replacement, organisation, law, infrastructure, and Earth-system
processes. This distinguishes the fixed individual-budget premise in \emph{Adaptability}'s Remark~\ref{AD-rem:race} from collective adaptation: partitioning, tools and organisational changes can change the effective bound. Neither a universal institutional rate nor its negation follows from biological development alone. Development and skill formation are responsive to social conditions
and cumulative investment \citep{nasem2019,heckman2006}; institutions can
coordinate under specified conditions, and can also fail through fragmented
authority, missing resources, and illegitimacy \citep{ostrom1990,ipcc2022wgii}.

Cost compression does not abolish physical conditions. IEA analysis reports
both possible operational gains from AI and rebound risks, and projects growing
data-centre electricity demand in its scenarios \citep{iea2025}. Its 2026
update stresses that per-task efficiency can fall while total demand rises with
volume and more intensive workloads \citep{iea2026}. Every proposed floor must
therefore report the energy, material, water, land, grid, labour, and local
burden boundary.

A realised GE graph substrate would enter the ledger only if an independently
reproducible implementation showed a lower $C^{\mathrm{aud}}$ at the same
pre-registered task, reliability, access, and externality conditions. Until
then, it is an unmeasured scenario: it may lower a technical component, leave
the total unchanged, or introduce new costs. It cannot be used to establish the
transition premise or accelerate a forecast.

\section{Falsifiers and conclusion}\label{sec:falsifiers}

The cost-floor thesis fails for a named case if: (i) performance-normalised,
boundary-harmonised lower envelope does not decline on the fixed outcome scale; (ii) the apparent saving is a transfer to an omitted payer, population,
time, or environmental sink; (iii) verified adaptation, access, and repair do
not improve despite the saving; or (iv) a comparable configuration provides the
outcome at lower total cost. The economic reorganisation hypothesis fails
locally if measured access, income security, agency, demand, and legitimacy
remain adequate under the declared task transformation. The transition
proposition does not apply unless $G(R_0)$ and the viability classification are
made explicit and independently contestable.

The evidence supports neither an inevitable collapse nor a guaranteed
post-scarcity successor. It supports a more disciplined question. Rapid falls
in selected information and coordination costs, together with longer-run
productivity and automation changes, can place pressure on arrangements built
around cost, work, and ownership. Whether that pressure becomes backlog,
adaptation, terminal discontinuity, or a viable successor regime is a matter
for declared measurements, not a title-level prophecy.

\begin{thebibliography}{99}
\sloppy
\bibitem[CMA(2024)]{cma2024} Competition and Markets Authority. (2024). \emph{AI Foundation Models: initial review}. \href{https://www.gov.uk/cma-cases/ai-foundation-models-initial-review}{Case page}.
\bibitem[Heckman(2006)]{heckman2006} Heckman, J. J. (2006). Skill formation and the economics of investing in disadvantaged children. \emph{Science}, 312(5782), 1900--1902. \href{https://doi.org/10.1126/science.1128898}{DOI}.
\bibitem[IEA(2025)]{iea2025} International Energy Agency. (2025). \emph{Energy and AI}. \href{https://www.iea.org/reports/energy-and-ai/executive-summary}{Executive summary}.
\bibitem[IEA(2026)]{iea2026} International Energy Agency. (2026). \emph{Key Questions on Energy and AI}. \href{https://www.iea.org/reports/key-questions-on-energy-and-ai/executive-summary}{Executive summary}.
\bibitem[ILO(2023)]{ilo2023} International Labour Organization. (2023). \emph{Generative AI and Jobs: A Global Analysis of Potential Effects on Job Quantity and Quality}, Working Paper 96. \href{https://www.ilo.org/sites/default/files/wcmsp5/groups/public/---dgreports/---inst/documents/publication/wcms_890761.pdf}{Working-paper PDF}.
\bibitem[ILO(2025)]{ilo2025} International Labour Organization. (2025). \emph{Policy Measures to Address Inequalities and Increase the Labour Income Share}. \href{https://researchrepository.ilo.org/view/pdfCoverPage?download=true&filePid=13138777410002676&instCode=41ILO_INST}{Publication PDF}.
\bibitem[IPCC(2022)]{ipcc2022wgii} IPCC. (2022). \emph{Climate Change 2022: Impacts, Adaptation and Vulnerability}. \href{https://www.ipcc.ch/report/ar6/wg2/}{Report page}.
\bibitem[NASEM(2019)]{nasem2019} National Academies of Sciences, Engineering, and Medicine. (2019). \emph{The Promise of Adolescence}. National Academies Press. \href{https://doi.org/10.17226/25388}{DOI}.
\bibitem[OECD(2021)]{oecd2021} OECD. (2021). \emph{Trade in the Time of Parcels}, OECD Trade Policy Papers 249. \href{https://www.oecd.org/content/dam/oecd/en/publications/reports/2021/05/trade-in-the-time-of-parcels_acf1943f/0faac348-en.pdf}{PDF}.
\bibitem[OECD(2025)]{oecd2025} OECD. (2025). \emph{Strengthening Supply Chains through Efficiency, Resilience, AI and Environmental Performance}. \href{https://www.oecd.org/en/publications/strengthening-supply-chains-through-efficiency-resilience-ai-and-environmental-performance_2a495d96-en.html}{Report page}.
\bibitem[Ostrom(1990)]{ostrom1990} Ostrom, E. (1990). \emph{Governing the Commons}. Cambridge University Press.
\bibitem[Our World in Data(2026)]{owidag2026} Our World in Data. (2026). \emph{In the Past, Most People Worked in Agriculture; in Today's Rich Countries, Only a Small Share Do}. Historical sources include Broadberry and Gardner (2013). \href{https://ourworldindata.org/data-insights/in-the-past-most-people-worked-in-agriculture-in-todays-rich-countries-only-a-small-share-do}{Data insight}.
\bibitem[Rissanen(1978)]{rissanen1978} Rissanen, J. (1978). Modeling by shortest data description. \emph{Automatica}, 14(5), 465--471. \href{https://doi.org/10.1016/0005-1098(78)90005-5}{DOI}.
\bibitem[Stanford HAI(2025)]{stanford2025} Stanford Institute for Human-Centered AI. (2025). \emph{AI Index Report 2025}. \href{https://hai.stanford.edu/assets/files/hai_ai_index_report_2025.pdf}{Report PDF}.
\bibitem[USDA ERS(2026)]{usda2026} U.S. Department of Agriculture, Economic Research Service. (2026). \emph{Productivity Growth in U.S. Agriculture}. Data updated June 2026. \href{https://www.ers.usda.gov/data-products/agricultural-productivity-in-the-united-states/productivity-growth-in-us-agriculture}{Data page}.
\fussy
\end{thebibliography}

\end{document}
